% Technical report for the team: the evidence that Ripple can
% replace C++ SCRIMMAGE. The academic paper is main.tex; it cites this report.
\documentclass[conference]{IEEEtran}
\usepackage{times}
\usepackage[numbers]{natbib}
\usepackage{booktabs}
\usepackage{tabularx}
\usepackage{array}
% Ragged-right paragraph column for narrow table cells.
\newcolumntype{P}[1]{>{\raggedright\arraybackslash}p{#1}}
\usepackage[bookmarks=true]{hyperref}

\newcommand{\stub}[1]{\par\noindent\textit{TODO: #1}\par}
% Marks results that do not yet come from a release run of the suite
% (reference/suite.py).
\newcommand{\prov}{\textbf{[Provisional]}~}

\begin{document}

\title{Ripple: Verifying a Rust Replacement for C++ SCRIMMAGE}
\author{Author names and affiliations to be added}
\maketitle

\begin{abstract}
This report gives the evidence for replacing C++ SCRIMMAGE with Ripple, a
Rust implementation that keeps SCRIMMAGE's mission
files and plugin design.
Run on the same missions, the two agree to floating-point rounding on 14
hand-written missions and 100 generated scenarios, comparing positions,
beliefs, sensor readings, plugin outputs, and every message delivery, and the
comparison catches all 12 planted defects. Every remaining difference is a C++
defect that Rust fixes or a documented choice, including a substep-order
regression in C++ that output comparison could not reveal. An experiment takes
less work to build and is checked before it runs; Rust is 1.04 to 1.6 times
faster on one core and uses 2 to 6 times less memory. Two extensions keep the
comparison passing: an estimation stage that separates truth, measurements, and
belief, and the removal of two scaling limits inherited from SCRIMMAGE. With
4{,}096 agents, a step with all-agent messages takes 9.4~ms against
378~ms in C++, and an optional spatial grid makes the step with collision
checks 2.4 times faster than scanning every agent, with byte-identical output. Every check is rerun from a clean checkout by one command.
\end{abstract}

\IEEEpeerreviewmaketitle

\section{Introduction}

SCRIMMAGE provides a plugin-based environment for multi-agent robotics
experiments \cite{demarco2018}. Ripple reimplements it in Rust so that
new work, higher-fidelity vehicles, heterogeneous fleets, and larger teams,
can build on a simulator that is deterministic at any worker count, checks
missions before running them, and is simpler to extend. Replacing a simulator
a team relies on needs evidence that the replacement behaves the same. This
report collects that evidence: what Ripple is
(Section~\ref{sec:baseline}), how closely it matches C++
(Section~\ref{sec:parity}), what it costs to build and run an experiment
(Sections~\ref{sec:workflow} and~\ref{sec:performance}), and two extensions that
keep the comparison passing (Sections~\ref{sec:estimation}
and~\ref{sec:scaling}). The details are in the appendices.

\section{A Validated Baseline}
\label{sec:baseline}
\label{sec:upgrades}

This section describes Ripple and summarizes three checks: that it
behaves like SCRIMMAGE, that an experiment is easier to build in it, and what
it costs to run. The behavior checks were repeated after adding the
estimation stage (Section~\ref{sec:estimation}), and the cost measurements
include it. The details are in the appendices.

\subsection{Ripple}

Each entity owns its autonomy, controllers, motion model, and sensors.
Interactions, networks, and metrics operate on the world. An autonomy chooses
a desired command; controllers turn it into motion-model inputs; the motion
model advances the state. Sensors then produce observations, interactions
apply effects such as collisions, networks deliver messages, and metrics
update scores. The engine finishes each phase before starting the next, even
when entities run on different CPU workers.

User plugins live in a workspace crate and register with the same
\texttt{ripple} executable as the stock models. A plugin owns its parameters
and state; the engine handles scheduling and message delivery. XML and YAML
missions load into one typed scenario; YAML adds shared entity templates and
parameter sweeps, and the simulator checks the scenario and plugin parameters
before running. A run saves trajectories, events, scores, and a manifest, with
Rerun for recorded visualization. Sweeps write one result row per case and can
be split into local processes or Slurm array shards, although cluster
submission has so far been tested with a mock scheduler.

\subsection{Behavior}
\label{sec:parity}

Does Ripple behave like SCRIMMAGE? We ran the same mission files in both
and compared not only where every agent went, but also what each agent
believed, what its sensors reported, what each plugin commanded, and when every
message arrived (Appendix~\ref{app:parity}).
\begin{itemize}
\item \emph{They agree.} All 14 hand-written missions and 100 randomly
generated scenarios match. Positions agree to under a nanometre, usually to
about $10^{-14}$~m, which is floating-point rounding.
\item \emph{The comparison is strict enough to catch real mistakes.} To test
it, we introduced 12 small, realistic errors into a copy of the Rust code, such
as a controller gain 1\% too high or a message arriving one step late. The
comparison flagged every one.
\item \emph{Long runs mostly hold.} Eight of ten runs of up to 1{,}000~s
match throughout. The other two end about 0.2~mm and $8\times10^{-9}$~m apart;
perturbation experiments suggest numerical sensitivity.
\item \emph{Every other difference is a C++ bug that Rust fixes or a
deliberate choice,} such as C++ scoring late-spawning teams as NaN; all are
listed in Table~\ref{tab:differences}.
\end{itemize}
This holds for the models and configurations we tested
(Appendix~\ref{app:coverage}), not the whole C++ plugin catalog.

\subsection{A Regression in the Reference}
\label{sec:substeps}

Matching a reference also means inheriting its mistakes. Reading the C++
history turned up one that output comparison could not, because both
simulators did the same thing.

\paragraph{The original design}
A mission can split each tick into shorter substeps, so that vehicle physics
runs at a finer time step than decisions and sensing. From its first public
commit in 2017 until December 2025, SCRIMMAGE ran each substep as the
controllers followed by the motion model. This is a cascade: a slow outer loop
(autonomy) sets goals, and a fast inner loop (controller and physics) tracks
them, with the controller seeing the state the previous substep wrote. JSBSim
likewise runs its flight-control system before each frame's dynamics
\cite{berndt2004jsbsim}, and Isaac Lab runs actuator models on every physics
step while holding a slower policy's action \cite{mittal2025isaaclab}.

\paragraph{What changed and why it matters}
The commit adding GPU support ran every controller substep before any motion,
so that one GPU call could advance all substeps together; it gives no reason
for the new order. With more than one substep, every controller call in a tick
then sees the same state and the same goal, and the motion model flies only the
last call's command. The controller gets no feedback within the tick, and a PID
controller's derivative term, the change in error divided by the time step, is
zero on every call whose setpoint and measurement are unchanged, so it loses its
damping. With one substep
the two orders are identical, which is why most missions were unaffected and
why comparison against the current C++ could not reveal it.

\paragraph{The fix}
The C++ comparison branch restores the interleaved order, and Ripple
follows it: controllers and motion alternate at each substep, while autonomy,
sensing, communication, and interactions keep the outer tick. Ripple
controllers read belief, so when own-state belief comes from the perfect
measurement it is refreshed after each motion substep, as the C++ controllers
read truth unless a sensor provides belief. Sensor-based estimates keep their
own timing, so more substeps do not mean fresher observations, or finer
collision checks. With one substep the two orders are the same, and the
unmodified and fixed C++ builds write byte-identical output on every such
mission that uses none of the other corrected plugins. The two substep missions
agree with the fixed C++ to $4.8\times10^{-14}$~m. Appendix~\ref{app:substeps} gives
the commits and the verification. The fix changes only the schedule; the
numerical approximations inside the retained motion models remain.

\subsection{Building an Experiment}
\label{sec:workflow}

We built the same small experiment in each simulator, following each one's own
documented workflow. One aircraft flies straight at 18~m/s; another chases it
with a new FollowNearest autonomy plugin. The question is how fast the chaser
must be to catch it, so we swept its speed over 18, 22, 26, and 30~m/s.
Table~\ref{tab:workflow} compares the steps.

\begin{table}[t]
\centering
\small
\caption{The FollowNearest speed sweep built in each simulator.}
\label{tab:workflow}
\begin{tabular}{@{}p{0.16\linewidth}p{0.37\linewidth}p{0.37\linewidth}@{}}
\toprule
Step & SCRIMMAGE (C++) & Ripple \\
\midrule
Mission & XML; each entity lists its own controller, motion model, and spread & YAML; entities share an \texttt{aircraft} template \\
Mistyped key & Ignored; runs with defaults & Rejected before the run, naming the expected key \\
Plugin & Generate a 41-file CMake project and 4 plugin files; fix the template by hand & One file in the starter crate; register in 3 lines \\
Build & CMake; set plugin and library paths & \texttt{cargo build} \\
Sweep & Mark values as \texttt{\$\{speed=30\}} in the XML; batch scripts broken, so write a loop & 9-line sweep file naming any key; \texttt{ripple sweep} \\
Cluster & \texttt{qsub} helper removed & Same sweep file via \texttt{bulk\_run.py} \\
Results & One \texttt{summary.csv} per log folder & One \texttt{results.jsonl} row per case \\
\bottomrule
\end{tabular}
\end{table}

Both give the same answer: the chaser catches the target at 26 and 30~m/s, at
45.3~s and 32.0~s, and not at 18 or 22~m/s. The difference is in the work
around the model. The mission files are the same length, but the YAML one
states shared setup once, and its sweep needs no edits to the mission. Checking
matters most: with the end time mistyped as \texttt{ending} and the speed as
\texttt{initial\_sped}, C++ ran without a warning for its default 50~s at the
plugin's default 24.3~m/s, and the chaser never caught the target. SCRIMMAGE's
own test missions \texttt{straight\_cpu\_mul.xml} and
\texttt{straight\_gpu\_mul.xml}
carry such a typo (\texttt{motion\_multipler}), so they never run substeps. In
YAML, the same mistakes stop the run before it starts.

The plugin code is similar in size (54 lines of C++ against 82 of Rust, which
also declares its ports and checks its parameter). On the C++ side, the
generated plugin does not compile until its placeholders are fixed, and neither
batch script runs, because both need a Python package that is no longer in the
branch. Appendix~\ref{app:workflow} gives the details, and \texttt{paper/workflow-example/} holds
both versions and the commands.

\subsection{Cost}
\label{sec:performance}

We ran the same five missions in both simulators on one laptop (an Apple M2
Pro): aircraft flying straight, the same with finer physics steps, aircraft
with noisy sensors, a stream of new aircraft that can collide, and a large
spread-out group checked for collisions every step. Each ran at several sizes,
up to 2{,}048 aircraft (Appendix~\ref{app:performance}).
\begin{itemize}
\item \emph{Rust is faster on one processor core,} by 1.04 to 1.6 times. We
compare speed only on the two missions where both simulators produce identical
results.
\item \emph{Rust uses much less memory,} 2.1 to 6.3 times less on every mission
and size.
\item \emph{Collision checking limits large teams, and spatial queries
remove the limit.} C++ compares every pair of aircraft, twice, so its cost
grows with the square of the team size: doubling from 1{,}024 to 2{,}048
aircraft made its collision mission 3.8 to 8.8 times slower. With
the optional grid of Section~\ref{sec:spatial}, a step with 4{,}096 spread-out
aircraft takes 3.8~ms, against 9.0~ms when every agent is scanned,
with identical output.
\item \emph{Message routing limits C++.} Deciding which messages reach which
plugins takes C++ one check per publisher--subscriber pair, so a sensor on
every agent that every agent listens to costs $n^2$ checks per tick. With
4{,}096 aircraft each running an own-state sensor, a C++ step takes
378~ms. Rust routes by index (Section~\ref{sec:messages}): the same step
takes 9.4~ms, with the same deliveries.
\item \emph{Extra cores help only when there is enough work.} Rust on eight
cores ran the sensor mission 3.3 times faster than on one, but the simplest
missions ran slower: handing out work to the cores every step costs more than
the work itself.
\item \emph{C++'s multi-core mode is unreliable.} It froze in 9 of 219 runs,
gave different answers to the same run when sensors were involved, and was
usually slower than running C++ on one core. We compare against single-core
C++.
\end{itemize}
We measured both in a Linux container and directly on macOS; the findings hold
in both.

\section{Separating Measurement from Belief}
\label{sec:estimation}

The first extension changes how a robot comes to know its own state and its
surroundings. In SCRIMMAGE a sensor can write the robot's belief directly, any
sensor can overwrite it, and a planner may read belief, subscribe to a sensor's
message itself, or fall back to truth when no estimate exists. Every
experiment that involves imperfect information has to wire these pieces by
hand. Ripple now has an explicit estimation stage: sensors produce
measurements, estimators own belief, and autonomy and controllers read only
that belief. Table~\ref{tab:estimation} shows what this changes for a
developer.

\begin{table*}[t]
\centering
\small
\caption{Common tasks before and after the estimation stage.}
\label{tab:estimation}
\begin{tabularx}{\textwidth}{@{}P{0.22\textwidth}XX@{}}
\toprule
Task & SCRIMMAGE & With the estimation stage \\
\midrule
Choose which sensor feeds the robot's state & Implicit: any sensor can overwrite belief, and the last writer in a step wins & Named in the mission and checked before the run \\
Give a planner its state & The planner reads belief, subscribes to a sensor's topic itself, or falls back to truth & The planner reads belief, always \\
Swap the localization method & Edit every planner that selected its own source & Change the localizer; planners are untouched \\
Configure error & Sensor noise, motion disturbance, and wiring are set in separate places & Each source of error is set once, at its owner \\
Keep decisions from using hidden truth & Not enforced: decisions can read the true state of every entity & Decisions receive belief only \\
Measure how wrong the robot was & Worked out by hand after the run & Truth and belief are recorded separately; recording measurements is future work \\
\bottomrule
\end{tabularx}
\end{table*}

\paragraph{Two estimators, one belief}
Belief has two parts, each with one owner. \emph{Localization} answers
``where am I?'': it turns navigation measurements (GPS, IMU) and the commands
the robot issued into its estimated state. \emph{Tracking} answers ``what is
around me?'': it turns contact measurements (radar, camera) into the contacts
the robot knows about, using the pose localization just produced. Each can be
replaced without touching the other or any planner. Sensors and evaluation
keep access to physical truth; estimators and decisions do not.

\paragraph{What changes in a tick}
The stage runs after sensors and message delivery and before metrics,
localization first, and keeps the existing order of every other phase. With a
0.1~s timestep:
\begin{enumerate}
\item \emph{Decide at time 0.0.} Autonomy and controllers read the belief for
time 0.0 and produce the inputs the motion model will use.
\item \emph{Advance to 0.1.} The motion model advances physical truth; sensors
sample the result, interactions apply their effects, and networks deliver the
messages due.
\item \emph{Form belief for 0.1.} Localization, then tracking, update belief
from the completed interval and the information now available.
\item \emph{Decide at time 0.1,} using that belief.
\end{enumerate}
Estimation adds computation to each tick, not simulated delay; modeling
processing delay would be a separate, explicit choice.

\paragraph{Pass-through first, prediction next}
The first implementation copies measurements into belief. With perfect, current
measurements, which simple missions use by default, belief equals truth at
each decision for every entity already observed (an entity that has just
appeared is known from the next measurement). Existing missions are therefore
unaffected: the C++ comparison, which compares every entity's belief as well as
its state, passes on every curated mission (Section~\ref{sec:parity}). The
stage runs inside the existing per-entity delivery step, so it adds work but no
extra hand-off between threads; the timings of Section~\ref{sec:performance}
include it. A predictive
localizer then occupies the same stage: it advances the previous belief using
the commands the robot issued and its own motion model, and corrects that
prediction with measurements. Because the physical response, the estimator's
assumed response, and measurement error are configured separately, an
experiment can vary each on its own. In the worked example
(\texttt{crates/starter}), a robot commanded at 10~m/s achieves 8~m/s and gets
an exact position fix every 5~s. With a localizer that assumes 10~m/s, the
position error grows by 2~m per second on straight segments, and the largest
error sampled at the end of a step is 9.8~m (the prediction each fix replaces
is 10~m off); with one that assumes 8~m/s, the estimated
position matches the true one. The starter crate's tests check these numbers.
Acting on the wrong belief, the robot turns early at every corner
of its square route and strays up to 22~m from a robot with perfect
localization.
% The timing, equality, startup, and removal rules are in the book
% (architecture/tick.md and information-flow.md); open items in docs/TODO.md.

\section{Removing Inherited Scaling Limits}
\label{sec:scaling}

Two costs in SCRIMMAGE's design grow with the square of the team size: finding
which agents are near each other, and deciding which messages reach which
plugins. Both are hidden inside the engine, and both come to dominate a step
as teams grow, in C++ as in Ripple, which kept the same design. They are
properties of the original design, not porting errors. We remove each inside
the engine, without changing any plugin's results, so experiments with larger
teams become practical.

\subsection{Finding Nearby Agents}
\label{sec:spatial}

Many models ask the same question: which agents are within some distance of a
point? Collision detection asks it with a few metres, a short-range radio with
a kilometre, a range-limited sensor with its detection range. SCRIMMAGE's
SimpleCollision answers by comparing every ordered pair of agents, so each pair
twice; its R-tree is used only to place agents at spawn time. Its collision
mission became 3.8 to 8.8 times slower from 1{,}024 to 2{,}048
aircraft (Appendix~\ref{app:performance}). Ripple moves this question into the
engine, so each plugin asks it once and the engine decides how to answer. With
4{,}096 spread-out agents, a whole step takes 9.0~ms when the engine
scans every agent and 3.8~ms with the optional grid, with identical output
(Table~\ref{tab:spatial}).

\paragraph{One question, two answers}
World plugins (sensors, networks, and metrics) call
\texttt{within(position, range)} on the true contacts. It returns every
contact whose distance from \texttt{position} is less than \texttt{range},
in the contacts' order. It does not filter: the caller itself, and contacts
already removed, are included, and each plugin applies its own rules.
Interactions, which change entities, ask the same question by entity index,
and can also ask for every close pair, \texttt{pairs\_within(range)}, in the
order a loop over all pairs would visit them.
The mission chooses how the engine answers, once for the whole run:
\begin{itemize}
\item \emph{Scan} (the default) tests every contact. It is the reference
answer, and it keeps the C++ timing comparison like-for-like.
\item \emph{Grid} (\texttt{spatial\_lookup: \{method: grid, cell\_width\_m:
$w$\}}) tests only contacts in nearby cells.
\end{itemize}
Plugin code is the same either way; no plugin knows which method answered.
XML missions always use the scan.

\paragraph{How the grid works}
After all motion substeps, the engine gives each agent an integer cell,
$(\lfloor x/w \rfloor, \lfloor y/w \rfloor, \lfloor z/w \rfloor)$, and adds
it to a hash map from its vertical $(x, y)$ column to the agents in that
column. Only occupied columns appear, so memory is one short entry per agent,
independent of the size of the map. A query with
range $r$ covers the cells overlapping a box of half-width $r$ around its
position, padded by a relative $10^{-9}$ so that rounding cannot exclude a
neighbour (typically nine columns when $r \le w$). The cells supply only
\emph{candidates}.
Every candidate then gets the scan's exact test, $\|p - q\| < r$
computed the same way, and the survivors are returned in entity order. The
grid can therefore only drop agents the exact test would reject, and never
changes the answer or its order; it changes only how many agents are tested.
If any position is not finite or too large for an integer cell, the tick uses
the scan. A query also scans when looking up its columns would cost more than
testing every agent; we measured one column lookup at about 30 distance tests.
Keying by column keeps lookups few when agents are spread out horizontally,
but a query tests every agent in its columns at any height, so many agents
stacked above one spot would bring back the every-pair cost.

\paragraph{When the grid is built}
After motion, before sensors, which share it without locks because it is built
before the parallel sensor phase. An interaction may move entities, so the
engine checks the index against current positions before each interaction and
again before networks and metrics, and rebuilds it if anything moved; the check
is one pass over the agents. Every query therefore sees current positions, and
the scan and the grid always answer from the same ones. SCRIMMAGE instead
builds its R-tree when entities are generated, before motion, so a plugin that
queries it later in the tick compares a moved agent with where the others
were. The index holds positions only: whether an agent is still alive is read
live by the plugin, because one collision can disable later ones.

\paragraph{Collision detection as the first user}
SimpleCollision previously visited every pair $(i, j)$ with $j > i$ in entity
order, skipped pairs with a disabled agent, and disabled both agents of each
pair closer than the collision range. The order matters: in a chain of three
overlapping agents, the pair visited first decides which agent survives. The
rewritten plugin iterates over \texttt{pairs\_within}: for each agent $i$ in
order, the engine finds its neighbours, keeps those with $j > i$, and yields
them in ascending order, as the plugin consumes them. That is exactly the old
sequence with the pairs that could never collide removed, so events and
survivors are unchanged. With the scan, the loop is the old double loop over a
compact list of positions. Placement at spawn time, which uses an
inclusive range and an altitude-only option, still scans, so spawn retries and
random draws are unchanged too.

\paragraph{What this preserves}
The determinism check (\texttt{rust-determinism},
\texttt{paper/data/rust-determinism/}) runs every mission in the repository,
the timing workloads at both sizes, including spawning with collisions
(153 collisions at 512 agents), and three crowded missions (512
to 2{,}048 agents, 235 to 795 collisions) with the scan at one and
eight workers and with grids from 0.7 to 1{,}000~m: frames, summaries, and
events are byte-identical. That the scan still matches SCRIMMAGE's every-pair
loop is the C++ comparison's job: its collision missions pass, and the planted
bug that makes the range test inclusive is caught. Tests cover the three-agent
chain
with ids created out of order, an interaction that moves one agent onto another
before SimpleCollision (the collision is found in the same step), points
exactly at the range, on cell faces, and at negative coordinates, contacts held
in a different order from their ids, and non-finite positions and ranges. The
grid tests use enough agents that the grid, not its fallback, answers. In the
timing runs below, every variant with a collision check also wrote identical
output, up to 4{,}096 agents.
Decisions cannot use the grid: autonomy and controllers see only belief, and
can ask the same \texttt{within} question of their own known contacts, which
the engine answers by scanning that list.

\paragraph{Cost}
Table~\ref{tab:spatial} times perf.py's collision mission with the scan and
with grids, alternating variants run by run (\texttt{speed-spatial},
\texttt{paper/data/speed-spatial/}). With 4{,}096 spread-out agents, a whole step takes 9.0~ms with the scan
and 3.8~ms with 5~m cells, 2.4 times faster; with 2{,}048 agents the grid is
1.6 times faster, and with 1{,}024, 1.2 times. Subtracting a run without
SimpleCollision estimates the check's own cost at 4{,}096 agents: about 6.0~ms
per step with the scan and 0.8~ms with the grid. It is only an estimate,
because that run removes no agents and may place spawns differently. Larger
cells were slightly faster (3.5~ms per step with 50~m cells), because a lookup
costs more than the few extra distance tests an occupied cell adds. The grid
can be slower than the scan: with 256 agents its queries fall back to the scan,
and building it makes the step 9\% slower (0.23 against 0.21~ms); in the
crowded 1{,}024-agent case the step is 10\% slower; and with a 50~m range and
5~m cells, 6\% slower, while 50~m cells there were 18\% faster than the scan.

\IfFileExists{generated/spatial.tex}{% Generated by paper/data/make_tables.py; do not edit.
\begin{table*}[t]
\centering
\caption{Stepping time in milliseconds per step with SimpleCollision, median of
5 runs after a warm-up, alternating variants, one worker, Darwin-arm64-10cpu-96207172.
Every variant of a case except the one without a collision check wrote identical
frames and events.}
\label{tab:spatial}
\footnotesize
\begin{tabular}{lrrrrrrrl}
\toprule
Case & Agents & Range (m) & Collisions & No collision check & Every pair & Scan & Grid 5\,m & Best other grid \\
\midrule
sparse & 256 & 5 & 0 & 0.19 & 0.22 & 0.21 & 0.23 & -- \\
sparse & 1{,}024 & 5 & 1 & 0.77 & 1.43 & 1.17 & 0.98 & -- \\
sparse & 2{,}048 & 5 & 7 & 1.67 & 5.21 & 3.26 & 2.11 & -- \\
sparse & 4{,}096 & 5 & 43 & 4.14 & 19.07 & 10.19 & 5.00 & 4.74 (50\,m) \\
clustered & 1{,}024 & 5 & 453 & -- & 0.27 & 0.27 & 0.28 & -- \\
clustered & 4{,}096 & 5 & 1930 & -- & 3.10 & 3.05 & 3.06 & -- \\
large range & 4{,}096 & 50 & 1421 & -- & 4.46 & 3.89 & 3.98 & 3.61 (50\,m) \\
\bottomrule
\end{tabular}
\par\smallskip Every pair: SimpleCollision's previous loop, at revision
b280a82. No collision check: the same mission without
SimpleCollision; the difference only estimates the check's cost, because that run removes
no agents and may place spawns differently. Sparse: aircraft spread over about a
kilometre at one altitude. Clustered: a 20\,m spread, so most agents collide early.
\end{table*}
}{}

\subsection{Delivering Messages}
\label{sec:messages}

Each tick, a network decides which published messages reach which subscribers.
SCRIMMAGE asks the network one question per link: can this publisher reach this
subscriber, and does the message get through? C++ keeps each topic's publishers
and subscribers in lists, so a message is checked against every subscriber to
its topic, and LocalNetwork, which delivers only within one agent, answers each
check by comparing the two plugins' parent entities. When every agent publishes
on a topic that every agent subscribes to, that is $n^2$ checks per tick to
deliver $n$ messages. The own-state sensor NoisyState does exactly this: every
agent publishes its reading on LocalNetwork, and every agent's Straight
autonomy subscribes to it. (In Ripple the reading reaches belief through
the estimation stage, and Straight discards the message; NoisyState keeps
publishing it only so that delivered messages still match C++. It is a
convenient built-in source of all-agent traffic, which auctions, shared maps,
and position broadcasts also produce.)

\paragraph{Routing by index}
Each routing phase, Ripple indexes the mailboxes subscribed to each channel
published that phase, in mailbox order, with the message type each expects, so
a message visits only its subscribers. A network that delivers only within the
sender's own agent says so by calling \texttt{route\_within\_entity} instead
of \texttt{route}; LocalNetwork does. The index is then kept per agent, so each
message sees only its own agent's subscribers, and a message from a world
plugin, which has no agent, sees none. The index is built only in phases that
have messages, so a network with nothing to route costs nothing extra, and
queued deliveries find their receivers in one map. Every link a network decides
is decided with the same inputs and in the same order as in C++, so deliveries
and random draws match. What the index cannot remove are links that exist: a
topic every agent publishes and subscribes to on a network that can reach every
agent (GlobalNetwork, or SphereNetwork in range) still has $n^2$ links, and
each must be decided.

\paragraph{What this preserves}
The C++ comparison traces every delivery, with its network, topic, sender, and
receiver, in order, and requires them to match C++ exactly: 6{,}132
deliveries over the curated and generated missions, including SphereNetwork's,
whose loss draws depend on the order in which links are decided. The planted
bug that lets LocalNetwork deliver across agents, a call to \texttt{route}
instead of \texttt{route\_within\_entity}, is caught by the trace comparison.

\paragraph{Cost}
Table~\ref{tab:messages} times perf.py's aircraft with and without NoisyState
in native C++ and in Rust (\texttt{speed-messages},
\texttt{paper/data/speed-messages/}). The difference between the two cases is everything the
sensor adds: its noise, readings, estimation, and messages. In C++ it grows with
the square of the team: 0.46, 6.2, 94, and 374~ms per step at 128, 512,
2{,}048, and 4{,}096 agents, against 3.7~ms for the rest of a 4{,}096-agent
step. A profile of the 2{,}048-agent case (saved beside the timings) puts
almost all of it in routing, in \texttt{LocalNetwork::is\_reachable} and the
\texttt{EntityPlugin::parent} calls it makes. In Ripple the sensor adds 0.16,
0.71, 3.0, and 6.5~ms per step, roughly linear in the team size, and no
routing function stands out in its profile: at 4{,}096 agents a whole step
takes 9.4~ms against 378~ms in C++, 40 times faster. Steps without messages
take 2.95~ms against 3.74~ms.

\IfFileExists{generated/messages.tex}{% Generated by paper/data/make_tables.py; do not edit.
\begin{table*}[t]
\centering
\caption{Milliseconds per step for perf.py's aircraft, without and with a NoisyState
sensor on every aircraft (whose messages every aircraft's Straight subscribes to): native
C++, Rust before routing by index, and after. Whole process minus a one-step run, median of
3 runs after a warm-up, alternating variants, one thread, Darwin-arm64-10cpu-96207172.
The two Rust builds wrote identical output.}
\label{tab:messages}
\footnotesize
\begin{tabular}{rrrrrrr}
\toprule
 & \multicolumn{3}{c}{Without messages} & \multicolumn{3}{c}{With NoisyState} \\
Agents & C++ & Rust before & Rust after & C++ & Rust before & Rust after \\
\midrule
128 & 0.16 & 0.10 & 0.10 & 0.62 & 0.55 & 0.27 \\
512 & 0.48 & 0.37 & 0.36 & 6.78 & 6.47 & 1.08 \\
2{,}048 & 1.67 & 1.39 & 1.36 & 95.92 & 94.88 & 4.38 \\
4{,}096 & 3.71 & 2.87 & 2.92 & 388.37 & 387.83 & 9.72 \\
\bottomrule
\end{tabular}
\end{table*}
}{}

\section{Limits}
\label{sec:discussion}


The C++ defects that Rust corrects are listed in Table~\ref{tab:differences},
and the regression Rust inherited in Section~\ref{sec:substeps}. The limits of
the parity claim are:
\begin{itemize}
\item \emph{Scope.} Parity covers the tested configurations of the retained
models, not the full C++ catalog.
\item \emph{Message contents.} Only lifecycle, collision, and sensor messages
are compared by content; other messages, such as auction bids, by delivery only.
\item \emph{Stochastic behavior.} Sensor noise is tested by seven marginal
distributions in one NoisyContacts fixture, not for temporal or cross-axis
dependence, and NoisyState noise is not tested by distribution.
\item \emph{Rounding.} Results agree to rounding, not bit for bit: C++ and Rust
call different math libraries and group some operations differently.
\end{itemize}

\section{Conclusion}

Ripple reproduces C++ SCRIMMAGE on every tested mission to
floating-point rounding, fixes the C++ defects it found on both sides, costs
less to build experiments in and to run, and removes two scaling limits without
changing any output. Within the scope of Appendix~\ref{app:coverage}, it can
replace the C++ simulator for the team's work.

\section*{Generative AI Disclosure}

OpenAI Codex assisted with the initial outline and reference-file
organization. Claude (Anthropic) assisted with implementation, running the
verification checks, and drafting the validation, agreement, and appendix text
from retained data. \stub{Update this statement to reflect actual assistance and
author review, following the selected venue's policy.}

\appendices

\section{Parity Details}
\label{app:parity}

\paragraph{Approach}
Both simulators run the same mission files. C++ comes from the current
\texttt{Ubuntu-24.04} branch, built headless in its own Docker image.

Trajectories alone can hide differences in what agents believe, what their
plugins command, and what messages they send. We therefore added an optional
trace to a separate C++ branch. It records each message delivery (tick,
sender, and receiver), the entities named in lifecycle and collision events,
published sensor readings, each agent's belief, and the commands each plugin
passes to the next. Other message contents, such as auction bids, are not
compared. Rust writes the same trace. Instrumentation that changed C++ would
invalidate the
comparison, so every run is repeated on the unmodified build: frames and
summaries must be byte-identical.

Randomized runs only compare if both simulators draw the same random numbers.
The C++ standard leaves its default generator up to each compiler's library, so
the C++ branch can switch spawning to the same fully specified generator Rust
uses. Rust's sensor noise uses its own random streams by design, so noise is
compared by distribution instead of sample by sample.

A comparison that cannot fail proves nothing. Beyond fourteen hand-written
missions, we ran 100 randomly generated scenarios and ten long ones of up to
1{,}000~s, planted twelve plausible porting bugs to confirm the comparison
catches them, and tested the noise distributions. Appendix~\ref{app:method}
gives the details.

\paragraph{Results}
Table~\ref{tab:parity} summarizes the evidence. All curated missions and
generated scenarios pass, and every planted bug is detected. Two
message-delivery bugs (a message one tick late, local messages crossing
entities) leave trajectories unchanged; only the summary and trace checks catch
them. In passing runs, trajectories agree to $2.4\times10^{-10}$~m at worst and
usually to $10^{-14}$~m or better. No difference was detected in seven
NoisyContacts marginal error distributions at the corrected significance
threshold (smallest $p = 0.031$ against 0.0071).

Eight of the ten long runs pass. The other two end about 0.2~mm and
$8\times10^{-9}$~m apart; perturbation experiments suggest numerical
sensitivity (Appendix~\ref{app:results}).

In the tested configurations of the retained models,
Ripple agrees with SCRIMMAGE in trajectories, the compared sensor and
command values, and message-delivery timing, to floating-point rounding.
Summaries match or are validated against the documented scoring differences.
The exceptions are listed in Table~\ref{tab:differences}.

These results come from one release run of the checking suite from a clean
checkout (Appendix~\ref{app:reproduction}); \texttt{paper/data/release.json}
records the hash of the source that ran and the C++ commits, \texttt{4bdf41fb0}
(\texttt{Ubuntu-24.04}) and \texttt{1c8595b91} (\texttt{benchmarking-edits}).

% Generated by paper/data/make_tables.py; do not edit.
\begin{table}[t]
\centering
\caption{Parity evidence. Appendix~\ref{app:results} has the full results.}
\label{tab:parity}
\small
\begin{tabularx}{\columnwidth}{P{0.34\columnwidth}>{\raggedright\arraybackslash}X}
\toprule
Curated missions & 13 of 13 pass \\
Generated scenarios & 100 of 100 pass \\
Long runs, up to 1,000 s & 8 of 10 pass; the others end $2.0\times10^{-4}$ m, $3.7\times10^{-8}$ m apart \\
Planted bugs & 12 of 12 detected \\
Sensor noise & no difference detected in 7 of 7 marginal distributions \\
Largest position difference, passing runs & $1.8\times10^{-10}$ m \\
Message deliveries compared & 126,199; 0 groups differ \\
\bottomrule
\end{tabularx}
\end{table}


\paragraph{What differs}
Apart from the long-run divergence above, every remaining difference
is a C++ defect that Rust corrects or a deliberate choice
(Table~\ref{tab:differences}). Examples are C++ scoring teams that spawn
late as NaN, and entities destroyed before the first step still running for a
tick. None is hidden. Where Rust's summary differs, entity and collision counts
still match exactly, and Rust's own messages reproduce the C++ summary under the
C++ scoring rule and the Rust summary under Rust's documented rule.

\section{Comparison Method}
\label{app:method}

\paragraph{Builds}
A checker builds each C++ branch from a \texttt{git archive} inside the
unmodified upstream dependency image (\texttt{linux/amd64}, GCC~13.3), headless
and single-threaded. \texttt{Ubuntu-24.04} is the unmodified reference;
\texttt{benchmarking-edits} adds opt-in comparison support and the substep
correction (Appendix~\ref{app:substeps}). A later
commit (\texttt{5777aeea1}) also makes AircraftPIDController use its given
timestep and lowers its default altitude derivative gain, matching Rust. The
timestep change is a fix; the lower gain is
a tuning choice, checked on one aircraft after one upset
(\texttt{paper/data/rk4-loads}, part C). A third commit
(\texttt{1c8595b91}) is the only change to model equations: FixedWing6DOF and
Multirotor computed forces and moments from the start-of-step state at every
RK4 stage, which made the step first order; both implementations now use each
stage's state (\texttt{paper/data/rk4-loads}). The release run reported here
includes all three. Two command drivers exist in both implementations:
\texttt{MotorSpeeds}, because the upstream multirotor controller no longer
matches the multirotor model, and \texttt{ConstantVelocity}, because no retained
upstream autonomy commands a velocity. Rust runs at 1, 2, and 8 workers, with
recording, and with tracing; all must be byte-identical.

\paragraph{Trace}
Set with \texttt{SCRIMMAGE\_TRACE} in C++ and \texttt{-{}-trace} in Rust. It
records every message delivery (tick, network, topic, sender, receiver, and the
entities in lifecycle and collision events), every published sensor state and
covariance, each entity's belief when its frame is logged, and every output the
next plugin reads. Each mission runs on the unmodified build, the instrumented
build, and the instrumented build with tracing; all three must write
byte-identical frames and summaries.

\paragraph{Criteria}
Frames must contain the same entities at every tick. States, beliefs, and
sensor readings must agree within $10^{-8}$~m, $10^{-8}$~m/s, $10^{-9}$~rad, and
$10^{-9}$~rad/s; outputs within $10^{-8}$; covariances within $10^{-9}$;
summaries within $10^{-6}$. Deliveries must match exactly, in order. The
position limit was set at 27 times the largest difference then measured in a
passing run ($3.7\times10^{-10}$~m); the largest in this release is
$2.4\times10^{-10}$~m.
They replace the original 0.01~m and $10^{-3}$, which let a planted
sub-micrometre bug through.

\paragraph{Random draws}
The C++ standard leaves the default random engine and the normal
distribution's algorithm to each standard library. The reference build's
library (GCC's) and the one Rust reproduces (LLVM's) differ in both. With
\texttt{SCRIMMAGE\_LIBCXX\_SPAWN\_RANDOM}, the C++ branch draws spawns from the
standard-defined \texttt{std::minstd\_rand} with the same normal algorithm as
Rust. Without it, the randomized mission differs by up to 8.42~m; with it, it
agrees to $1.4\times10^{-14}$~m. Sensor noise is compared with
fixed-bias, zero-noise fixtures and, separately, by distribution.

\paragraph{Long runs}
Ten generated scenarios run for 1{,}000~s (fixed-wing: 100~s at 0.01~s steps).
A run agrees if every check passes over its whole length. For a run that
diverges, a control changes Rust's initial roll, commanded speeds, and turning
radii by one unit in the last place, scales each controller loop's proportional
gain by $1+10^{-12}$, and measures when Rust diverges from itself. One-ulp
parameter changes alone perturb the state far less than the rounding difference
between C++ and Rust, and do not reach a controller loop directly.

In \texttt{generated-000}, one aircraft settles into a sustained pitch
oscillation at about 20~s. Its altitude loop never levels off, and its commanded
pitch rate alternates between about $+1$ and $-3.7$~rad/s every half second.
Any difference in that loop then grows about tenfold per second until it levels
off near 0.1~mm. C++ and Rust exceed the tolerance after 28.8~s and end
$2.0\times10^{-4}$~m apart, and the control does the same after 26.8~s,
ending $2.1\times10^{-4}$~m away (Table~\ref{tab:long}). This suggests
numerical sensitivity; it does not rule out a small difference between the two
implementations.

\texttt{generated-007} flies four FixedWing6DOF aircraft with
AircraftPIDController. C++ and Rust exceed the tolerance after 76.4~s and end
$8.4\times10^{-9}$~m apart; the control diverges much earlier, after 14.4~s,
and ends $7.9\times10^{-7}$~m away (Table~\ref{tab:long}). Again this
suggests numerical sensitivity rather than a porting error, with the same
caveat. In \texttt{generated-009} the control diverges after 33~s while C++
and Rust agree throughout.

\paragraph{Summaries and one exclusion}
A differing summary passes only if entity and collision counts match exactly,
and the Rust run's metrics messages reproduce the C++ summary under the C++
scoring rule and the Rust summary under Rust's documented rule
(\texttt{reference/metrics\_model.py}). For entities C++ destroys before
the first step, their first-step trace records are excluded and counted, because
C++ still runs them for that step. No other difference is excluded.

\section{Full Results}
\label{app:results}

Table~\ref{tab:agreement} gives each curated mission, Table~\ref{tab:generated}
the generated scenarios, Table~\ref{tab:sensitivity} the planted bugs,
Table~\ref{tab:noise} the noise distributions, Table~\ref{tab:long} the long
runs, and Table~\ref{tab:differences} every known difference. Retained data is
under \texttt{paper/data/}.

% Generated by paper/data/make_tables.py; do not edit.
\begin{table*}[t]
\centering
\caption{C++/Rust behavioral agreement against the instrumentation branch.
Maximum absolute errors over all compared entity states: position in m,
velocity in m/s, attitude as the rotation angle between orientations in rad.
Trace columns count compared message deliveries, published sensor states,
and plugin output values; every delivery matched in sender, receiver, tick,
and order. In every case the C++ comparison build wrote byte-identical frames
and summaries with and without tracing, and matched the unmodified build
except where substeps differ by design (Section~\ref{sec:substeps}); Rust
outputs were byte-identical at 1, 2, and 8 workers, with recording, and with
tracing.}
\label{tab:agreement}
\scriptsize
\setlength{\tabcolsep}{2.5pt}
\begin{tabular}{lP{0.14\textwidth}rllllrrrll}
\toprule
& & & \multicolumn{3}{c}{Frames} & \multicolumn{4}{c}{Trace} & \\
\cmidrule(lr){4-6}\cmidrule(lr){7-10}
Mission & Exercises & States & Position & Velocity & Attitude & Deliv. & Payload & Outputs & Output err. & Result \\
\midrule
\texttt{straight-no-gui} & SimpleAircraft, PID, collision, removal & 976 & 0 & 0 & 0 & 5 & 0 & 5856 & 0 & match \\
\texttt{straight\_cpu} & SimpleAircraft, PID, long horizon & 2001 & 0 & 0 & 0 & 1 & 0 & 12000 & 0 & match \\
\texttt{straight\_cpu\_mul} & as above; upstream multiplier typo & 2001 & 0 & 0 & 0 & 1 & 0 & 12000 & 0 & match \\
\texttt{aircraft-substeps-spawning} & substeps, rates, randomized spawning & 90 & $2.2\times10^{-14}$ & $1.9\times10^{-14}$ & $7.8\times10^{-15}$ & 4 & 0 & 516 & $9.1\times10^{-14}$ & match \\
\texttt{aircraft-substeps-scheduled} & substeps, rates, two spawn schedules & 170 & $4.8\times10^{-14}$ & $2.6\times10^{-14}$ & $7.8\times10^{-15}$ & 7 & 0 & 978 & $1.8\times10^{-13}$ & match\textsuperscript{a} \\
\texttt{noisy-state-bias} & NoisyState bias, belief feedback & 31 & 0 & $3.6\times10^{-15}$ & 0 & 31 & 30 & 180 & $3.6\times10^{-15}$ & match \\
\texttt{networks-local-global} & Boundary, GroundCollision, networks & 607 & $5.9\times10^{-14}$ & $4.5\times10^{-14}$ & $5.1\times10^{-15}$ & 613 & 605 & 3630 & $5.4\times10^{-14}$ & match\textsuperscript{a} \\
\texttt{fixed-wing-6dof} & FixedWing6DOF, AircraftPID, wind & 2002 & $6.9\times10^{-13}$ & $8.8\times10^{-14}$ & $4.1\times10^{-15}$ & 2 & 0 & 14000 & $5.0\times10^{-13}$ & match \\
\texttt{fixed-wing-6dof-rascal} & FixedWing6DOF Rascal, upset, glide & 18003 & $2.3\times10^{-13}$ & $7.3\times10^{-14}$ & $4.0\times10^{-15}$ & 3 & 0 & 126000 & $8.1\times10^{-14}$ & match \\
\texttt{multirotor} & Multirotor (4 and 6 rotors) & 16004 & $5.0\times10^{-14}$ & $4.8\times10^{-13}$ & $3.6\times10^{-14}$ & 4 & 0 & 72000 & 0 & match \\
\texttt{single-integrator} & SingleIntegrator command paths & 204 & 0 & 0 & 0 & 4 & 0 & 1200 & 0 & match \\
\texttt{noisy-contacts-bias} & NoisyContacts bias, attitude, covariance & 93 & $7.1\times10^{-14}$ & $1.6\times10^{-14}$ & $4.4\times10^{-15}$ & 93 & 180 & 540 & $5.5\times10^{-14}$ & match \\
\texttt{sphere-network-auction} & SphereNetwork reach, AuctionAssign & 426 & 0 & 0 & 0 & 22 & 0 & 2520 & 0 & match \\
\texttt{collision-boundary} & SimpleCollision range boundary & 108 & 0 & 0 & 0 & 14 & 0 & 660 & 0 & match\textsuperscript{a} \\
\bottomrule
\end{tabular}
\par\smallskip\footnotesize
\textsuperscript{a}The summary differs from C++ and is reproduced exactly by
applying the C++ scoring rule to the Rust run's metrics messages
(Table~\ref{tab:differences}).
\end{table*}

% Generated by paper/data/make_tables.py; do not edit.
\begin{table}[t]
\centering
\caption{Differential testing over 100 generated scenarios
(\texttt{generate\_scenarios.py}, seed 1), with the same checks as
Table~\ref{tab:agreement}.}
\label{tab:generated}
\footnotesize
\begin{tabularx}{\columnwidth}{P{0.42\columnwidth}>{\raggedright\arraybackslash}X}
\toprule
Scenarios passing & 100 of 100 \\
Compared entity states & 67,973 \\
Compared deliveries / outputs & 5,328 / 437,571 \\
Max position / velocity error & $3.1\times10^{-11}$ m / $2.5\times10^{-11}$ m/s \\
Max attitude / output error & $1.0\times10^{-12}$ rad / $2.1\times10^{-11}$ \\
Summaries identical / explained & 51 / 49 \\
Pre-start removal records excluded & 308 \\
Scenarios with collisions & 51 \\
Scenarios with ground & 37 \\
Scenarios with boundary & 23 \\
Scenarios with scheduled spawns & 61 \\
Scenarios with NoisyState & 13 \\
Scenarios with NoisyContacts & 16 \\
Scenarios with fixed-wing & 21 \\
Scenarios with point-mass & 63 \\
\bottomrule
\end{tabularx}
\end{table}

% Generated by paper/data/make_tables.py; do not edit.
\begin{table}[t]
\centering
\caption{Planted defects: 12 of 12 detected. Each
is a one-line change applied to a copy of the Rust source and compared
with the saved C++ runs of Table~\ref{tab:agreement}. Missions counts the
matrix missions that failed; the layers are those that reported it.}
\label{tab:sensitivity}
\footnotesize
\begin{tabularx}{\columnwidth}{>{\raggedright\arraybackslash}Xrl}
\toprule
Defect & Missions & Detected by \\
\midrule
controller substeps receive the full step dt & 2 & frames, trace \\
frame timestamp labels the step start, not t + dt & 14 & frames, summary \\
spawn coordinates skip the C++ six-decimal round trip & 1 & frames, trace \\
network delivers one tick late & 14 & summary, trace \\
SphereNetwork range becomes inclusive & 1 & trace \\
SimpleCollision range becomes inclusive & 1 & frames, summary, trace \\
sensor attitude errors applied in yaw-pitch-roll order & 2 & frames, trace \\
NoisyContacts covariance is identity, not 5I & 1 & trace \\
LocalNetwork delivers across entities & 2 & trace \\
SingleIntegrator yaw swaps atan2 arguments & 2 & frames, trace \\
heading PID proportional gain 1\% high & 5 & frames, trace \\
SimpleAircraft yaw rate 0.1\% high & 5 & frames, trace \\
\bottomrule
\end{tabularx}
\end{table}

% Generated by paper/data/make_tables.py; do not edit.
\begin{table}[t]
\centering
\caption{NoisyContacts error distributions, C++ / Rust, from
3,600 samples per quantity and simulator. Two-sample Kolmogorov-Smirnov
test; the Bonferroni threshold is 0.0071. Position
and velocity errors in m and m/s; attitude as the error rotation angle.}
\label{tab:noise}
\footnotesize
\setlength{\tabcolsep}{3pt}
\begin{tabular}{lllll}
\toprule
Quantity & Mean & Std. dev. & $D$ & $p$ \\
\midrule
pos\_x & 1.121 / 1.017 & 4.970 / 4.980 & 0.0203 & 0.446 \\
pos\_y & -2.017 / -1.850 & 5.012 / 4.927 & 0.0183 & 0.577 \\
pos\_z & 0.518 / 0.484 & 1.003 / 0.979 & 0.0264 & 0.161 \\
vel\_x & 0.111 / 0.096 & 4.999 / 4.965 & 0.0181 & 0.597 \\
vel\_y & 0.381 / 0.275 & 4.981 / 4.995 & 0.0339 & 0.031 \\
vel\_z & -0.021 / 0.019 & 0.989 / 0.994 & 0.0214 & 0.379 \\
attitude\_rad & 0.141 / 0.141 & 0.072 / 0.071 & 0.0169 & 0.676 \\
\bottomrule
\end{tabular}
\end{table}

% Generated by paper/data/make_tables.py; do not edit.
\begin{table}[t]
\centering
\caption{Long runs. When C++ and Rust first differ by more than the $10^{-8}$~m
tolerance, and, as a control, when Rust differs from itself after a
one-unit-in-the-last-place change to initial roll, commanded speed, and turning
radius, with each controller loop's proportional gain scaled by $1+10^{-12}$.
Final: largest C++/Rust position difference at the end.}
\label{tab:long}
\footnotesize
\setlength{\tabcolsep}{2pt}
\begin{tabular}{lllll}
\toprule
Scenario & Length & C++ vs Rust & Rust vs nudged & Final \\
\midrule
\texttt{generated-000} & 1000 s & 28.8 s & 26.8 s & $2.0\times10^{-4}$ \\
\texttt{generated-001} & 1000 s & never & never & $8.1\times10^{-12}$ \\
\texttt{generated-002} & 1000 s & never & never & $3.6\times10^{-12}$ \\
\texttt{generated-003} & 1000 s & never & never & $7.3\times10^{-12}$ \\
\texttt{generated-004} & 1000 s & never & never & $3.6\times10^{-12}$ \\
\texttt{generated-005} & 1000 s & never & never & 0 \\
\texttt{generated-006} & 1000 s & never & never & $9.0\times10^{-12}$ \\
\texttt{generated-007} & 100 s & 76.4 s & 14.4 s & $8.4\times10^{-9}$ \\
\texttt{generated-008} & 100 s & never & never & $7.0\times10^{-11}$ \\
\texttt{generated-009} & 1000 s & never & 33.0 s & $5.5\times10^{-12}$ \\
\bottomrule
\end{tabular}
\end{table}


\begin{table*}[t]
\centering
\caption{Known C++/Rust differences and how the comparison treats them.}
\label{tab:differences}
\small
\begin{tabularx}{\textwidth}{P{0.18\textwidth}>{\raggedright\arraybackslash}XP{0.22\textwidth}}
\toprule
Difference & Cause & Treatment \\
\midrule
Survivor flight time &
C++ publishes \texttt{EntityPresentAtEnd} after its last network phase, so
metrics credit survivors only until the latest removal. Rust credits them until
the end. &
Summary must equal the C++ rule applied to Rust's messages. \\
Teams first spawned after $t=0$ &
C++ scores such a team without a normalization window: infinite
\texttt{flight\_time\_norm}, NaN \texttt{score}. &
As above. \\
Entities destroyed before the first step &
C++ keeps them for one step: they still run their plugins, are reported removed
twice, and score as surviving. Rust removes them once. &
First-step trace records excluded and counted; summary as above. \\
Randomized spawning &
C++ leaves its default random generator to the standard library; the reference
build's differs from the one Rust uses. &
C++ branch uses the same generator as Rust. \\
Sensor noise samples &
Rust gives each sensor its own random stream; C++ shares one. &
Compared by distribution. \\
SphereNetwork positions &
C++ queries a range tree built before motion; Rust uses current post-motion
positions. &
No compared delivery differed. \\
FixedWing6DOF inertia &
C++'s default silently overrides a supplied SI inertia matrix; Rust accepts
either and rejects both together. &
Fixtures use the default. \\
AircraftPIDController with substeps &
C++ uses the full step as the controller timestep during substeps; Rust uses
the substep. &
FixedWing6DOF scenarios run without substeps; \texttt{5777aeea1} fixes the
comparison branch. \\
\bottomrule
\end{tabularx}
\end{table*}

\section{The Substep Correction}
\label{app:substeps}

In \texttt{SimControl::run\_entities}, SCRIMMAGE interleaved controller and
motion substeps from its first public commit (August 2017) until commit
\texttt{c0a9f42b2} (``added gpu support'', December 2025):
\begin{verbatim}
autonomy                      once per tick
for each substep:
    controllers               read the state motion just wrote
    motion model
sensors, interactions, ...    once per tick
\end{verbatim}
That commit closed the controller loop before a new GPU motion call, which
advances all substeps at once, and passed each controller call the full tick
length; a later commit restored the substep length but kept the split. A
comment written afterward still describes the interleaved design. GPU motion
models need OpenCL and appear in two demonstration missions.

The derivative term vanishes when a controller call sees the same setpoint and
the same measurement as the previous one: in the split order, autonomy runs
once per tick and nothing moves between controller calls. A nested controller
can still shift an inner setpoint between calls (an outer loop's integral keeps
growing), so the inner derivative is small rather than exactly zero.

The comparison branch (\texttt{7129e47f5}) restores the interleaved loop and
disables GPU acceleration in CMake, which also disables its GPU network plugin;
the GPU sources remain but are not built. \texttt{reference\_check.py} requires
the unmodified and fixed C++ to write identical output only when a mission has
one substep (and none of the other corrected plugins); in this release they do
on all nine curated missions that qualify. Against the fixed C++, the two
substep missions agree to $4.8\times10^{-14}$~m with matching traces, and the
unmodified C++ differs from the fixed build on both.

\section{Performance Details}
\label{app:performance}

Tables~\ref{tab:benchmark} and~\ref{tab:benchmark-native} time five workloads
(straight flight, ten controller and motion substeps, noisy state and contact
sensors, scheduled spawning with collisions, and collision checks among aircraft
spread over about a kilometre, so few collide and the population stays nearly
constant) on an Apple M2 Pro, first with
both simulators in one Linux container (\texttt{reference/benchmark.py}) and
then on bare-metal macOS, where C++ is built with Apple clang
against the library versions Ubuntu 24.04 ships
(\texttt{paper/data/speed-cpp-container/} and
\texttt{paper/data/speed-cpp-native/}). Rust output matches single-threaded
C++ in the first two workloads. The other three depend on random draws that the
C++ build makes differently without the comparison spawn stream, so their ratios
are not speed claims.
C++ is the comparison branch of Appendix~\ref{app:substeps}, so both simulators
run substeps in the same order.

Where outputs match, Rust with one thread is 1.07 to 1.6 times faster than C++
in the container and 1.04 to 1.4 times faster on bare metal, and it uses 2.1 to
6.3 times less memory on every workload. These timings include the estimation
stage.

C++ collision checking compares every pair of aircraft, twice, so its cost
grows with the square of the number of aircraft; Rust uses the spatial queries
of Section~\ref{sec:spatial}, by default the scan. From 1{,}024 to 2{,}048
aircraft, Rust's collision workload takes 2.7 to 2.8 times longer and C++'s 3.8
to 8.8 times longer. Eight workers barely help, because the check runs in one
sequential phase.

Parallel speedup depends on each phase carrying enough work: every phase hands
entities to a thread pool and joins it. On bare metal, eight Rust workers beat
one with substeps at 1{,}024 aircraft (1.2 times), with sensors at 64 and 128
aircraft (3.3 times at 128), and with collisions at 2{,}048 aircraft (1.1
times), but not on light workloads or with 32 sensing aircraft; in the
container, only with 128 sensing aircraft (1.6 times). Rust runs single-worker
phases on the calling thread, without the pool.

C++'s own \texttt{multi\_threaded} mode is not reliable. It deadlocked in 9 of 219
runs, counting startup runs (each killed after 120~s and retried), its output
with sensors differs between identical runs, and it is slower than
single-threaded C++ except with 64 or more sensing aircraft on bare metal and
128 in the container. Its entity
worker waits on a condition variable without a predicate,
so a wake-up can be lost; we did not modify it, and single-threaded C++ is the
baseline.

\IfFileExists{generated/benchmark.tex}{% Generated by paper/data/make_tables.py; do not edit.
\begin{table*}[t]
\centering
\caption{Stepping time in seconds (whole process minus a one-step run),
median of 3 runs after one warm-up, both simulators built for aarch64 in one Linux container.
C++~8 is C++'s own \texttt{multi\_threaded} mode.}
\label{tab:benchmark}
\footnotesize
\setlength{\tabcolsep}{4pt}
\begin{tabular}{lrrrrrrrrr}
\toprule
Workload & Agents & C++ 1 & C++ 8 & Rust 1 & Rust 8 & Ratio & C++ MiB & Rust MiB & Hangs \\
\midrule
\texttt{motion} & 64 & 0.161 & 1.872 & 0.132 & 3.852 & 1.22 & 14 & 4 & 2 \\
\texttt{motion} & 256 & 0.622 & 5.200 & 0.427 & 4.903 & 1.46 & 24 & 5 & 0 \\
\texttt{motion} & 1024 & 2.599 & 18.785 & 1.934 & 5.900 & 1.34 & 66 & 11 & 0 \\
\texttt{substeps} & 64 & 0.146 & 2.186 & 0.107 & 4.678 & 1.36 & 14 & 4 & 0 \\
\texttt{substeps} & 256 & 0.548 & 5.956 & 0.401 & 6.187 & 1.37 & 24 & 5 & 0 \\
\texttt{substeps} & 1024 & 2.576 & 21.334 & 1.554 & 5.351 & 1.66 & 65 & 11 & 0 \\
\texttt{sensing} & 32 & 0.271 & 0.630 & 0.251 & 1.028 & -- & 13 & 4 & 0 \\
\texttt{sensing} & 64 & 1.029 & 1.410 & 0.875 & 1.327 & -- & 17 & 6 & 0 \\
\texttt{sensing} & 128 & 5.120 & 4.241 & 3.124 & 1.952 & -- & 32 & 12 & 0 \\
\texttt{churn} & 64 & 0.190 & 1.679 & 0.114 & 3.492 & -- & 14 & 4 & 5 \\
\texttt{churn} & 256 & 0.565 & 3.357 & 0.242 & 4.149 & -- & 18 & 5 & 1 \\
\texttt{churn} & 1024 & 3.938 & 10.658 & 0.662 & 5.167 & -- & 35 & 7 & 1 \\
\texttt{collision} & 256 & 0.175 & 0.637 & 0.049 & 0.479 & -- & 24 & 5 & 0 \\
\texttt{collision} & 1024 & 2.607 & 4.616 & 0.232 & 0.719 & -- & 64 & 11 & 0 \\
\texttt{collision} & 2048 & 23.766 & 26.766 & 0.642 & 0.850 & -- & 118 & 21 & 0 \\
\bottomrule
\end{tabular}
\par\smallskip Ratio: C++ over Rust at one thread, shown only where their
outputs match; sensing, churn, and collision draw different random numbers, so the two
simulators do different work. MiB: peak resident memory.
Hangs: C++~8 runs killed after 120\,s and retried (see text).
\end{table*}
}{}
\IfFileExists{generated/benchmark-native.tex}{% Generated by paper/data/make_tables.py; do not edit.
\begin{table*}[t]
\centering
\caption{Stepping time in seconds (whole process minus a one-step run),
median of 3 runs after one warm-up, both simulators built natively on bare-metal macOS (arm64).
C++~8 is C++'s own \texttt{multi\_threaded} mode.}
\label{tab:benchmark-native}
\footnotesize
\setlength{\tabcolsep}{4pt}
\begin{tabular}{lrrrrrrrrr}
\toprule
Workload & Agents & C++ 1 & C++ 8 & Rust 1 & Rust 8 & Ratio & C++ MiB & Rust MiB & Hangs \\
\midrule
\texttt{motion} & 64 & 0.101 & 0.986 & 0.097 & 0.469 & 1.05 & 10 & 4 & 0 \\
\texttt{motion} & 256 & 0.390 & 2.564 & 0.359 & 0.805 & 1.09 & 19 & 6 & 0 \\
\texttt{motion} & 1024 & 1.605 & 8.999 & 1.375 & 1.542 & 1.17 & 54 & 12 & 0 \\
\texttt{substeps} & 64 & 0.127 & 1.230 & 0.094 & 0.517 & 1.35 & 10 & 4 & 0 \\
\texttt{substeps} & 256 & 0.506 & 3.438 & 0.358 & 0.742 & 1.41 & 19 & 6 & 0 \\
\texttt{substeps} & 1024 & 1.990 & 12.281 & 1.400 & 1.257 & 1.42 & 54 & 13 & 0 \\
\texttt{sensing} & 32 & 0.262 & 0.323 & 0.200 & 0.216 & -- & 10 & 5 & 0 \\
\texttt{sensing} & 64 & 0.931 & 0.689 & 0.682 & 0.392 & -- & 14 & 7 & 0 \\
\texttt{sensing} & 128 & 3.493 & 1.989 & 2.495 & 0.784 & -- & 28 & 13 & 0 \\
\texttt{churn} & 64 & 0.135 & 0.912 & 0.091 & 0.461 & -- & 10 & 4 & 0 \\
\texttt{churn} & 256 & 0.482 & 2.021 & 0.199 & 0.679 & -- & 14 & 6 & 4 \\
\texttt{churn} & 1024 & 3.111 & 6.450 & 0.600 & 1.108 & -- & 27 & 9 & 1 \\
\texttt{collision} & 256 & 0.177 & 0.399 & 0.047 & 0.096 & -- & 19 & 6 & 0 \\
\texttt{collision} & 1024 & 2.432 & 3.180 & 0.221 & 0.237 & -- & 53 & 12 & 0 \\
\texttt{collision} & 2048 & 9.304 & 10.288 & 0.591 & 0.560 & -- & 99 & 21 & 0 \\
\bottomrule
\end{tabular}
\par\smallskip Ratio: C++ over Rust at one thread, shown only where their
outputs match; sensing, churn, and collision draw different random numbers, so the two
simulators do different work. MiB: peak resident memory.
Hangs: C++~8 runs killed after 120\,s and retried (see text).
\end{table*}
}{}

\section{Coverage}
\label{app:coverage}

Compared: SimpleAircraft, FixedWing6DOF, Multirotor, and SingleIntegrator with
their controllers; Straight; AuctionAssign's message-delivery timing and
topology; NoisyState and
NoisyContacts; SimpleCollision, GroundCollision, and Boundary; Global, Local,
and SphereNetwork; and SimpleCollisionMetrics. Also compared: motion substeps,
update rates, scheduled and randomized spawning, and startup placement.
AuctionAssign's bids are random and its result message differs on purpose, so
its bids and winners are not compared. WaypointFollower, WaypointBroadcast,
NoisyPosition,
and the two command drivers are Rust-only.

\section{Rust Self-Consistency}
\label{app:self}

Rust results are byte-identical across 1, 2, and 8 workers, with and without
recording and tracing, and whether a scenario comes from XML, YAML, or Rust
code. Built for another operating system and processor (Linux/amd64 versus
macOS/arm64), events and summaries stay identical in all 56 runs checked (every
mission at 1 and 8 workers) and trajectories differ by at most
$2.4\times10^{-12}$~m, because each system's math
library rounds differently (\texttt{reference/platform\_check.py}). This is
exact because workers own separate entities, every phase finishes
before the next starts, and random draws are tied to entities and plugins
rather than workers.

\section{Checks}
\label{app:suites}
\label{app:refactors}

One runner, \texttt{reference/suite.py}, runs every check by a stable ID; each
has a purpose, a reference, and a pass criterion (Table~\ref{tab:checks}).
The quick set runs after a change; the release set adds the slower C++
studies, benchmarks, and focused studies. A check's evidence is retained in
\texttt{paper/data/<check>/} only when every selected check passed, and
\texttt{paper/data/release.json} records each outcome with the hash of the
source that ran. The \texttt{rust-tests} check runs \texttt{cargo test} (plugin
contracts, every curated mission at 1, 2, and 8 workers with and without
tracing, the estimation stage, XML/YAML pairs, the Rust scenario API, the
command-line runner and sweeps), \texttt{cargo fmt}, \texttt{cargo clippy},
and the unit tests of the checking tools and bulk-run scripts.

Speed regressions are judged against one accepted baseline, recorded with its
workloads, machine, and compiler; a measurement on another machine or
toolchain, or of changed workloads, is not compared and fails until a new
baseline is accepted by hand. Changes meant to preserve behavior are checked
with \texttt{reference/equivalence.py}, which runs every mission, the timing
workloads, and a message-heavy mission at 1, 2, and 8 workers with the trace
under two builds and requires byte-identical frames, events, summaries, and
traces.

\IfFileExists{generated/checks.tex}{% Generated by paper/data/make_tables.py; do not edit.
\begin{table*}[t]
\centering
\caption{The checks of \texttt{reference/suite.py} and their outcome in this
release (\texttt{paper/data/release.json}); they ran on 4 different sources.
}
\label{tab:checks}
\footnotesize
\begin{tabularx}{\textwidth}{lXP{0.22\textwidth}P{0.22\textwidth}}
\toprule
Check & What it checks & Against & Outcome \\
\midrule
\texttt{speed-regression} & Rust speed and memory have not regressed &
  reference/perf\_baseline.json, accepted on the same machine and toolchain &
  passed; no workload slower than the baseline \\
\texttt{rust-tests} & Rust unit, integration, and mission regression tests;
  formatting; lints; tests of these checking tools & expected values in the
  tests & passed \\
\texttt{rust-determinism} & worker counts, recording, tracing, and spatial
  lookup method change no output & the one-worker run & passed; 41 missions, 0
  differing \\
\texttt{cpp-agreement} & behaviour matches C++ SCRIMMAGE on the curated missions
  (reference/missions.txt) & C++ SCRIMMAGE in Docker: the unmodified and
  comparison branches & passed; 14 of 14 pass \\
\texttt{cpp-generated} & agreement over 100 random scenarios (seed 1) & C++
  SCRIMMAGE & passed; 100 of 100 pass \\
\texttt{cpp-long} & agreement over ten 1,000 s scenarios (seed 7), and whether
  late divergence is the port or the dynamics & C++ SCRIMMAGE; Rust after
  one-ulp nudges & passed; 8 of 10 agree throughout \\
\texttt{cpp-noise} & sensor noise distributions match where samples cannot & C++
  SCRIMMAGE NoisyContacts samples & passed; random payloads differ, as expected
  \\
\texttt{checker-sensitivity} & the C++ comparison catches plausible porting bugs
  & cpp-agreement's saved C++ runs & passed; 12 of 12 planted bugs caught \\
\texttt{rust-platform} & Linux/amd64 and native builds agree & the native build
  & passed; largest position difference 2.3678883310186693e-12 m \\
\texttt{speed-cpp-container} & C++ and Rust speed and memory in one Linux
  container & C++ SCRIMMAGE comparison branch, same container & passed \\
\texttt{speed-cpp-native} & C++ and Rust speed and memory on bare-metal macOS &
  native C++ SCRIMMAGE comparison branch & passed \\
\texttt{speed-spatial} & collision checking cost with the scan and with grids &
  the scan, which tests every contact & passed \\
\texttt{speed-messages} & message routing cost with all-agent traffic & native
  C++ SCRIMMAGE & passed \\
\texttt{rk4-loads} & what evaluating forces at every RK4 stage changes &
  unmodified C++ SCRIMMAGE, which holds forces through the step & passed \\
\texttt{multirotor-control} & MultirotorVelocityController's step responses on
  the Crazyflie, and convergence as the step shrinks & the commands; Förster's
  rotor drag; the 20-substep run & passed \\
\texttt{rotors-motor} & Multirotor's motor lag and rotor loads against RotorS,
  sample by sample & RotorS's motor-plugin source (ethz-asl/rotors\_simulator),
  compiled in Docker & passed \\
\texttt{crazyflow-crazyflie} & Multirotor's Crazyflie against Crazyflow's
  first-principles dynamics, same constants and motor commands & Crazyflow
  (learnsyslab/crazyflow, pinned), run with uv & passed \\
\texttt{stall-landing} & whether stall tables change ArduPlane's Rascal landing
  & the same mission without tables, and a repeat of it & passed \\
\bottomrule
\end{tabularx}
\end{table*}
}{}

\section{Workflow and Visual Checks}
\label{app:workflow}

Local bulk runs with 3 and 10 shards reproduced the eight-case waypoint sweep's
complete \texttt{results.jsonl}. Collection rejected missing and mismatched
shards. Slurm submission has been tested only against a mock scheduler; no
cluster campaign has run.

Visual review of \texttt{networks-local-global} (Rerun 0.36.2, 1920$\times$1200,
checkpoints at 0, 0.3, 0.6, 10, and 29.9~s) confirmed map framing, labels,
status, complete paths, and ground removal
(\texttt{paper/data/visual-qa/}). The speed plot of a nearly constant-speed
aircraft autoscaled into unreadable spikes, so the dashboard has not passed a
complete visual review. The waypoint demonstrations have not been reviewed with
screenshots.

\section{The Workflow Example}
\label{app:workflow}

The C++ side used the unmodified \texttt{Ubuntu-24.04} reference image.
\texttt{create-scrimmage-project.py} generated the project, and
\texttt{generate-plugin.sh autonomy FollowNearest} generated the header,
source, parameter XML, and CMake file. The templates on this branch were
reformatted to \texttt{(>{}>{}> PLUGIN\_NAME < < <)}, which the generator's
substitution no longer matches, so four placeholders in the source and three in
the header had to be fixed by hand. The template's example logic already
follows the nearest opponent, so the plugin body is unchanged from it. Running
the plugin requires adding the build and header folders to
\texttt{SCRIMMAGE\_PLUGIN\_PATH} and \texttt{LD\_LIBRARY\_PATH}.
\texttt{run\_experiments.py} and \texttt{scrimmage\_runner.sh} both import
\texttt{scrimmage.utils}, which was removed from the branch, and need
\texttt{pyDOE}, \texttt{pandas}, and \texttt{tqdm}; the \texttt{scrimmage}
command has no option to set a mission variable. The sweep is therefore a
shell loop that writes one mission per speed with \texttt{sed}.
\texttt{paper/workflow-example/run-cpp.sh} repeats all of this in one command.

The Rust side is the book's first-plugin tutorial: \texttt{follow\_nearest.rs}
in the starter crate, three lines in its \texttt{lib.rs}, a 48-line YAML
mission, and a 9-line sweep file. The same sweep file runs locally with
\texttt{ripple sweep} or as Slurm array shards with
\texttt{scripts/bulk\_run.py}. Line counts exclude blank and comment lines;
the C++ count covers the header and source, not the generated CMake file.
Both place the two aircraft at fixed positions, and all four outcomes and
both catch times agree.

\section{Reproduction}
\label{app:reproduction}

Every result in this report comes from one command, run in a fresh copy of the
repository in a differently named folder, with no earlier runs and no Git
history, beside the C++ checkout on both branches and with Docker running:
\begin{verbatim}
python3 reference/suite.py --release --retain
\end{verbatim}
It runs every check of Table~\ref{tab:checks}, replaces their evidence in
\texttt{paper/data/}, and regenerates the tables. \texttt{reference/README.md}
documents each check, the builds, tolerances, trace records, and how to add a
model; the compared missions are listed in \texttt{reference/missions.txt}.
The C++ branch \texttt{benchmarking-edits} holds the instrumentation as
reviewable commits on top of \texttt{Ubuntu-24.04}. Only the workflow example
runs separately (\texttt{paper/workflow-example/run-cpp.sh} and \texttt{ripple sweep}).

\bibliographystyle{plainnat}
\bibliography{../references}

\end{document}
